\documentclass[12pt]{article}
\usepackage[utf8]{inputenc}
\usepackage[T1]{fontenc}
\usepackage[margin=1in]{geometry}
\usepackage{setspace}
\usepackage{xcolor}
\usepackage[colorlinks=true, urlcolor=blue!60!black, linkcolor=blue!60!black]{hyperref}
\setlength{\parindent}{0em}
\setlength{\parskip}{0.6em}
\renewcommand{\baselinestretch}{1.0}
\begin{document}
\begin{center}
{\LARGE \textit{United States v.\ Skrmetti}, 605 U.S.\ 495 (2025)\par}
\vspace{0.5em}
{\large Case Brief --- Two-Page Edition\par}
\end{center}

\section{Memory Jogger}
Tennessee's ban on puberty blockers and hormones for minors' gender dysphoria classifies by age and medical use, so rational basis review applies (\href{https://www.courtlistener.com/opinion/10776831/united-states-v-skrmetti/}{CourtListener}).

\section{Detailed Case Facts}
Tennessee's SB1 bars giving minors puberty blockers or hormones to treat gender dysphoria or to live inconsistently with biological sex, while allowing the drugs for congenital defects, precocious puberty, disease, or injury. Transgender adolescents, parents, and a doctor sued; the United States intervened.

\section{Procedural History}
The district court enjoined the medication ban under intermediate scrutiny; the Sixth Circuit reversed under rational basis review.

\section{Judicial Votes}
Affirmed, 6--3. Roberts wrote for Thomas, Gorsuch, Kavanaugh, Barrett, and (in part) Alito; Thomas, Barrett, and Alito concurred; Sotomayor (with Jackson and, in part, Kagan) and Kagan dissented.

\section{Holding}
SB1 classifies by age and medical use, not sex or transgender status, and rationally protects minors amid medical uncertainty. The Court left open whether transgender status is a suspect class.

\section{Analysis of Opinions}
\textbf{Majority:} A treatment is drug plus diagnosis; SB1 forbids nothing to one sex that it permits the other, and under \textit{Geduldig} it does not single out transgender status. \textit{Bostock}'s but-for test fails because changing sex would not supply a qualifying diagnosis. \textbf{Barrett and Alito:} transgender status is not a suspect class. \textbf{Thomas:} courts should not defer to ``self-described experts.'' \textbf{Sotomayor:} SB1 classifies by sex on its face, as hirsutism and gynecomastia treatments show.

\section{Examples: Future Applications}
\textbf{Same-Side:} a Medicaid exclusion of gender-dysphoria treatments, likely a diagnosis-based line. \textbf{Opposite-Side:} a dress code barring clothing ``inconsistent with'' biological sex. \textbf{Fence-Sitter:} barring minors from changing birth-certificate sex markers.

\section{Critique}
Defenders praise leaving pediatric medicine to legislatures; Kate Redburn faults ``double deference'' in the medical context.

\section{Key Quotations}
\begin{quote}
``In the medical context, the mere use of sex-based language does not sweep a statute within the reach of heightened scrutiny.''
\end{quote}
\noindent --- Chief Justice Roberts, for the Court.

\end{document}
